\documentclass[12pt]{article}
\usepackage{amsmath}
\usepackage{amsthm}
\usepackage{amssymb}

\newtheorem{theorem}{Theorem}

\title{Introduction to Galois Field}
\author{Thanh Vuong Dang}
\date{19/05/2026}

\begin{document}
	\maketitle
	Galois fields,  named after Evariste Galois also known as Finite Field,  is a mathematical concept in abstract algebra that deals with finite mathematical structures. It is a set of numbers that consists of a finite number of elements and has two operations, addition and multiplication, that follow specific rules. The rules for these operations ensure that the Galois Field remains closed, meaning that the result of any operation performed within the set will also be an element of the set.
	
	Galois Fields are useful in various fields, such as cryptography, coding theory, and error correction, due to their unique mathematical properties.
	
	The size of a Galois Field is represented by a prime number 'p', and it is denoted by GF(p), where p is a prime number.
	\text{Example} One classical example of a Galois Field is a field with 2 elements, denoted as $GF(2)$ with two math operations which are addition and multiplication.
	
	\textbf{Properties of Galois Field}
	\begin{itemize}
		\item[1] \textbf{Finite Size:} The most important property of a Galois Field is that it is finite. It has a specific number of elements, and it is not possible to add any more elements to it. The size of the field is represented by a prime number 'p'.
		\item[2] \textbf{Closure:} The Galois Field remains closed under both addition and multiplication operations, meaning that the result of any operation performed within the set will always be an element of the set.
		\item[3] \textbf{Commutative:} The Galois Field is commutative under both addition and multiplication operations. This means that the order of elements does not matter in performing operations. For example, a+b = b+a and ab = ba.
		\item[4] \textbf{Associative:} The Galois Field is associative under both addition and multiplication operations. This means that the grouping of elements in an operation does not matter. For example, (a+b) + c = a + (b+c) and (ab) * c = a * (bc).
		\item[5] \textbf{Distributive:} The Galois Field follows the distributive property. This means that multiplication distributes over addition, i.e., a * (b + c) = a * b + a * c.
		\item[6] \textbf{Identity Elements:} The Galois Field has two identity elements, 0 for addition and 1 for multiplication. Any element added with 0 is equal to the original element, and any element multiplied by 1 is equal to the original element.
		\item[7] \textbf{Inverse Elements:} Every element in the Galois Field has an inverse element under both addition and multiplication operations. The inverse element for addition is the negative of the original element, and the inverse element for multiplication is the reciprocal of the original element.
	\end{itemize}

\section{Algorithm of polynomial division over $GF(2)$}
    Given two polynomial $A(x)$ and $B(x)$ on $GF(2)$, we set $Q(x)$ is the quotient of division $A/B$ and $R(x)$ is the corresponding remainder. Initialize $Q(x) = 0$ and $R(x) = A(x)$. Let $M = deg(A)$. \\
    While $M \geq deg(B)$, do: 
    \begin{itemize}
        \item[1] Let $$t(x) = \epsilon x^{M-deg(B)},$$ 
        where $\epsilon = 1$ if $M = deg(R)$, otherwise $\epsilon = 0$.
        \item[2] Update quotient: 
        $$ Q(x) = Q(x) \oplus t(x).$$
        \item[3] Update remainder:
        $$R(x) = R(x) \oplus t(x)B(x).$$
        \item[4] Update iteration step: 
        $$M = M - 1.$$
    \end{itemize}
    Stop when $M < deg(B).$ \\
    Return $A(x) = Q(x)B(x) + R(X).$
\end{document}
